\section{Prototype Implementation}\label{sec:implementation}

This section describes how a functional prototype realizes the five-tier hybrid architecture presented in Section~\ref{sec:model}. The prototype validates the smart-contract coordination layer (Tier~4) and demonstrates end-to-end data flow across community capture (Tier~1), off-chain encrypted storage (Tier~5), selective anchoring, hierarchical access control, and provenance tracking. Edge/fog preprocessing (Tier~2) and district-level batching (Tier~3) are implemented as application-layer modules; Layer-2 rollup and zero-knowledge verification components are specified architecturally and remain targets for future engineering. Figure~\ref{fig:proto_arch} presents the prototype software architecture mapped to the five-tier topology.

\begin{figure}[!ht]
\centering
\fbox{\parbox{0.8\columnwidth}{\centering\vspace{2em}[Prototype Software Architecture]\\Figure placeholder\vspace{2em}}}
\caption{Prototype software architecture mapped to the five-tier hybrid topology. Solid modules are implemented; dashed modules are architecturally specified.}\label{fig:proto_arch}
\end{figure}

% ===== Mermaid Source =====
% Figure: Prototype Software Architecture
% Scientific purpose: Maps the implemented prototype modules to the five architectural
% tiers, distinguishing fully implemented components from architecturally specified ones.
%
% graph TB
%     subgraph T1["Tier 1 · Community Capture Module"]
%         UI["React/Next.js Client"]
%         ENC_C["Client-Side AES-256 Encryption"]
%         MM["MetaMask Wallet"]
%     end
%
%     subgraph T2["Tier 2 · Edge/Fog Module"]
%         VAL_M["Validation Middleware"]
%         BUF_M["Local Buffer Queue"]
%     end
%
%     subgraph T3["Tier 3 · District Coordination Module"]
%         BATCH["Batch Aggregator"]
%         SUP["Supervisory Review Interface"]
%     end
%
%     subgraph T4["Tier 4 · Smart Contract Layer"]
%         RF["ReportFactory"]
%         RC["ReportCoordinator"]
%         ACL_C["AccessControl"]
%         COMP["CompletenessVerifier"]
%     end
%
%     subgraph T5["Tier 5 · Off-Chain Storage"]
%         IPFS_N["IPFS Node (ipfs-http-client)"]
%         ENC_S["Encrypted Payload Store"]
%     end
%
%     UI -->|"encrypted payload"| ENC_C
%     ENC_C -->|"ciphertext"| VAL_M
%     VAL_M --> BUF_M
%     BUF_M -->|"validated events"| BATCH
%     BATCH -->|"anchor tx"| RF
%     RF --> RC
%     RF --> ACL_C
%     RF --> COMP
%     ENC_C -.->|"encrypted blob"| IPFS_N
%     IPFS_N --- ENC_S
%     RC -.->|"CID pointer"| IPFS_N
%     MM -->|"sign tx"| RF
%     SUP -->|"review/rework"| RC
%
%     style T1 fill:#f5f5f5,stroke:#333
%     style T2 fill:#e8e8e8,stroke:#333
%     style T3 fill:#dcdcdc,stroke:#333
%     style T4 fill:#d0d0d0,stroke:#333
%     style T5 fill:#c4c4c4,stroke:#333
% =========================

\subsection{Development Environment and Technology Stack}\label{subsec:devenv}

The prototype was implemented using the following technology stack: Solidity~0.8 for smart contract development targeting the Ethereum Virtual Machine; Node.js as the server-side runtime; React and Next.js for the front-end application layer with server-side rendering; Web3.js for blockchain interaction; \texttt{ipfs-http-client} for IPFS content-addressed storage; MetaMask for client-side transaction signing; Ganache~CLI for local blockchain emulation during development; Mocha for automated contract testing; and Semantic~UI for interface theming. Client-side encryption uses the Web Crypto API with AES-256-GCM. Smart contracts were compiled using \texttt{solc} and deployed via a custom deployment script using the Truffle~HD Wallet Provider for account authentication.

The project comprises six modules: (1)~Ethereum components (smart contracts, compilation, deployment, and migration scripts); (2)~IPFS integration (encryption, upload, retrieval, and content-hash management); (3)~React page components implementing Tier~1 capture interfaces; (4)~Next.js routing, server-side rendering, and API middleware implementing Tier~2 validation logic; (5)~supervisory and batch coordination interfaces implementing Tier~3 workflows; and (6)~test components for contract and integration verification.

\subsection{Tier-by-Tier Implementation}\label{subsec:tiers}

\textbf{Tier~1: Community Data Capture.} Community health workers and facility reporters interact through a React-based web client. Upon case report submission, the client encrypts the surveillance payload using AES-256-GCM before any network transmission. The encrypted ciphertext is uploaded to IPFS via \texttt{ipfs-http-client}, which returns a content identifier (CID). The client then initiates a blockchain transaction through MetaMask, passing only the CID, the payload hash, and permission metadata to the smart contract---no raw surveillance data leaves the client unencrypted, and no raw data reaches the blockchain. This design accommodates lightweight client devices: reporters require only a browser with MetaMask and connectivity sufficient to reach an IPFS gateway and an Ethereum node endpoint.

\textbf{Tier~2: Edge/Fog Preprocessing.} The Next.js API middleware layer implements preliminary validation, schema conformance checks, and local buffering. When a reporter submits an event, the middleware validates field completeness and data-type conformance before forwarding the encrypted payload for IPFS upload. During connectivity interruptions, submissions are buffered in a local queue and retransmitted upon reconnection. This preprocessing reduces invalid submissions reaching the blockchain layer and provides near-source error feedback to reporters. In the current prototype, edge/fog logic executes within the Next.js server; production deployment would distribute these functions to dedicated edge nodes at health facility sites.

\textbf{Tier~3: District Coordination.} District and regional coordinators access a supervisory interface that aggregates validated submissions from subordinate reporters. The coordination module supports batch review: coordinators inspect submission metadata (timestamps, reporter identifiers, CIDs, and completeness status) without accessing decrypted case data unless explicitly authorized. Upon batch approval, the coordinator triggers a single anchoring transaction that registers multiple CIDs and their associated metadata on-chain, reducing per-report transaction costs. Rework requests and supervisory feedback are recorded as on-chain provenance events linked to the originating report CID.

\textbf{Tier~4: Blockchain Coordination Layer.} The smart contract layer is described in Section~\ref{subsec:contracts}. On-chain state is strictly limited to: cryptographic hashes, IPFS CID pointers, access-control lists, provenance and consent records, completeness verification state, alert events, and audit trail entries. Raw surveillance payloads are never stored on-chain.

\textbf{Tier~5: Off-Chain Encrypted Storage.} Encrypted surveillance payloads are stored on IPFS using content addressing. Each payload is encrypted at the originating client (Tier~1) before upload; the resulting CID serves as both a storage pointer and an integrity anchor. The smart contract stores only the CID and a SHA-256 hash of the plaintext for independent integrity verification. Authorized retrievers obtain the CID from the smart contract, fetch the encrypted blob from IPFS, and decrypt locally using keys distributed through the permission management system. This separation ensures that neither the blockchain nor the IPFS layer alone exposes identifiable health data.

\subsection{Smart Contract Architecture}\label{subsec:contracts}

The prototype implements four interdependent smart contracts that realize the selective anchoring principle described in Section~\ref{subsec:dataflow}. Figure~\ref{fig:contracts} illustrates the contract interaction topology.

\begin{figure}[!ht]
\centering
\fbox{\parbox{0.8\columnwidth}{\centering\vspace{2em}[Smart Contract Interaction Topology]\\Figure placeholder\vspace{2em}}}
\caption{Smart contract interaction topology implementing selective anchoring, access control, and completeness verification.}\label{fig:contracts}
\end{figure}

% ===== Mermaid Source =====
% Figure: Smart Contract Interaction Topology
% Scientific purpose: Shows how the four smart contracts interact to implement
% selective anchoring, permission management, provenance tracking, and completeness
% verification without storing raw surveillance data on-chain.
%
% graph LR
%     subgraph Contracts["Smart Contract Layer (Tier 4)"]
%         RF["ReportFactory<br/>· createReport()<br/>· getReports()<br/>· anchorBatch()"]
%         RC["ReportCoordinator<br/>· anchorHash()<br/>· recordProvenance()<br/>· requestRework()<br/>· markReviewed()"]
%         ACL["AccessControl<br/>· grantAccess()<br/>· revokeAccess()<br/>· checkAccess()<br/>· setExpiry()"]
%         CV["CompletenessVerifier<br/>· registerExpected()<br/>· recordSubmission()<br/>· getCompleteness()"]
%     end
%
%     subgraph OnChain["On-Chain State"]
%         H["CID + Hash Registry"]
%         P["Permission & Consent State"]
%         PR["Provenance Trail"]
%         CS["Completeness Ratios"]
%         AE["Alert & Audit Events"]
%     end
%
%     RF -->|"deploy instance"| RC
%     RF -->|"register"| CV
%     RC --> H
%     RC --> PR
%     RC --> AE
%     ACL --> P
%     CV --> CS
%     CV -->|"gap alert"| AE
%
%     style Contracts fill:#e8e8e8,stroke:#333
%     style OnChain fill:#d0d0d0,stroke:#333
% =========================

\textbf{ReportFactory.} The factory contract serves as the system entry point. It deploys and manages \texttt{ReportCoordinator} instances for each reporting unit, maintains a registry of all deployed coordinator addresses, and supports batch anchoring of multiple submissions in a single transaction. The factory pattern ensures that each coordinator retains an unmodified copy of the contract source, preventing post-deployment tampering.

\textbf{ReportCoordinator.} Each coordinator instance manages the lifecycle of surveillance reports for a specific reporting unit. It stores only selective on-chain artifacts: the IPFS CID, a SHA-256 hash of the original plaintext for integrity verification, reporter address, timestamp, and tier-level metadata. The coordinator supports provenance operations including supervisory review confirmation, rework requests with descriptive annotations, and hierarchical feedback---all recorded as immutable on-chain events. Listing~\ref{lst:anchor} illustrates the selective anchoring interface.

\begin{figure}[!ht]
\begin{verbatim}
// SPDX-License-Identifier: MIT
pragma solidity ^0.8.0;

contract ReportCoordinator {
    struct Anchor {
        string  ipfsCid;
        bytes32 contentHash;
        address reporter;
        uint256 timestamp;
        uint8   tierLevel;
        bool    reviewed;
    }

    mapping(uint256 => Anchor) public anchors;
    uint256 public anchorCount;

    event ReportAnchored(
        uint256 indexed id,
        string  ipfsCid,
        bytes32 contentHash,
        address reporter,
        uint8   tierLevel
    );
    event ReportReviewed(
        uint256 indexed id,
        address reviewer
    );
    event ReworkRequested(
        uint256 indexed id,
        address supervisor,
        string  annotation
    );

    function anchorReport(
        string  calldata _cid,
        bytes32 _hash,
        uint8   _tier
    ) external returns (uint256) {
        uint256 id = anchorCount++;
        anchors[id] = Anchor(
            _cid, _hash, msg.sender,
            block.timestamp, _tier, false
        );
        emit ReportAnchored(
            id, _cid, _hash, msg.sender, _tier
        );
        return id;
    }

    function markReviewed(uint256 _id)
        external
    {
        anchors[_id].reviewed = true;
        emit ReportReviewed(_id, msg.sender);
    }

    function requestRework(
        uint256 _id,
        string calldata _note
    ) external {
        emit ReworkRequested(
            _id, msg.sender, _note
        );
    }
}
\end{verbatim}
\caption*{Listing~1: Selective anchoring and provenance interface (\texttt{ReportCoordinator}).}\label{lst:anchor}
\end{figure}

\textbf{AccessControl.} The access-control contract implements the dynamic permission model described in Section~\ref{subsec:privacy}. Permissions are scoped by tier level, reporter--supervisor relationship, and data category. The contract supports three permission operations: \texttt{grantAccess} assigns read permission to a specific address for a given report identifier; \texttt{revokeAccess} removes a previously granted permission; and \texttt{setExpiry} attaches a time-bound deadline after which the permission automatically becomes invalid. Permission state changes are emitted as on-chain events for auditability. Listing~\ref{lst:acl} illustrates the access-control interface.

\begin{figure}[!ht]
\begin{verbatim}
contract AccessControl {
    struct Permission {
        bool    granted;
        uint256 expiry;   // 0 = no expiry
        uint8   tierScope;
    }

    // reportId => grantee => Permission
    mapping(uint256 => mapping(
        address => Permission
    )) public permissions;

    event AccessGranted(
        uint256 indexed reportId,
        address grantee,
        uint256 expiry
    );
    event AccessRevoked(
        uint256 indexed reportId,
        address grantee
    );

    function grantAccess(
        uint256 _reportId,
        address _grantee,
        uint256 _expiry,
        uint8   _tierScope
    ) external {
        permissions[_reportId][_grantee] =
            Permission(true, _expiry, _tierScope);
        emit AccessGranted(
            _reportId, _grantee, _expiry
        );
    }

    function revokeAccess(
        uint256 _reportId,
        address _grantee
    ) external {
        permissions[_reportId][_grantee]
            .granted = false;
        emit AccessRevoked(_reportId, _grantee);
    }

    function checkAccess(
        uint256 _reportId,
        address _grantee
    ) external view returns (bool) {
        Permission memory p =
            permissions[_reportId][_grantee];
        if (!p.granted) return false;
        if (p.expiry != 0 &&
            block.timestamp > p.expiry)
            return false;
        return true;
    }
}
\end{verbatim}
\caption*{Listing~2: Dynamic access control with grant, revoke, and time-bound expiry (\texttt{AccessControl}).}\label{lst:acl}
\end{figure}

\textbf{CompletenessVerifier.} The completeness contract implements the event-driven verification mechanism described in Section~\ref{subsec:completeness}. For each reporting period, expected reporting units are registered. As anchoring events arrive, the contract increments a submission counter. The completeness ratio---received submissions divided by expected units---is queryable on-chain without accessing underlying payload data. When the ratio falls below a configurable threshold at period close, the contract emits a gap alert event. Listing~\ref{lst:completeness} illustrates the completeness verification interface.

\begin{figure}[!ht]
\begin{verbatim}
contract CompletenessVerifier {
    struct Period {
        uint256 expectedCount;
        uint256 receivedCount;
        bool    closed;
    }

    // periodId => Period
    mapping(uint256 => Period)
        public periods;
    // periodId => reporter => submitted
    mapping(uint256 => mapping(
        address => bool
    )) public submitted;

    event SubmissionRecorded(
        uint256 indexed periodId,
        address reporter
    );
    event GapAlert(
        uint256 indexed periodId,
        uint256 received,
        uint256 expected
    );

    function registerPeriod(
        uint256 _periodId,
        uint256 _expectedCount
    ) external {
        periods[_periodId] = Period(
            _expectedCount, 0, false
        );
    }

    function recordSubmission(
        uint256 _periodId
    ) external {
        require(
            !submitted[_periodId][msg.sender],
            "Already submitted"
        );
        submitted[_periodId][msg.sender] = true;
        periods[_periodId].receivedCount++;
        emit SubmissionRecorded(
            _periodId, msg.sender
        );
    }

    function getCompleteness(
        uint256 _periodId
    ) external view returns (
        uint256 received, uint256 expected
    ) {
        Period memory p = periods[_periodId];
        return (p.receivedCount, p.expectedCount);
    }
}
\end{verbatim}
\caption*{Listing~3: Completeness verification with gap alerting (\texttt{CompletenessVerifier}).}\label{lst:completeness}
\end{figure}

\subsection{IPFS Integration and Off-Chain Storage Workflow}\label{subsec:ipfs_impl}

The off-chain storage workflow implements the Tier~5 data layer. Figure~\ref{fig:ipfs_workflow} illustrates the IPFS anchoring lifecycle.

\begin{figure}[!ht]
\centering
\fbox{\parbox{0.8\columnwidth}{\centering\vspace{2em}[IPFS Anchoring Workflow]\\Figure placeholder\vspace{2em}}}
\caption{IPFS anchoring workflow: encryption, upload, CID return, on-chain anchoring, and authorized retrieval.}\label{fig:ipfs_workflow}
\end{figure}

% ===== Mermaid Source =====
% Figure: IPFS Anchoring Workflow
% Scientific purpose: Shows the complete lifecycle of a surveillance payload from
% client-side encryption through IPFS upload to on-chain CID anchoring and
% authorized retrieval with decryption.
%
% sequenceDiagram
%     participant Client as Reporter Client
%     participant Crypto as AES-256-GCM
%     participant IPFS as IPFS Gateway
%     participant SC as ReportCoordinator
%     participant ACL as AccessControl
%     participant Retriever as Authorized Retriever
%
%     Client->>Crypto: Encrypt payload
%     Crypto-->>Client: ciphertext + IV + tag
%     Client->>IPFS: Upload encrypted blob
%     IPFS-->>Client: Return CID
%     Client->>Client: Compute SHA-256(plaintext)
%     Client->>SC: anchorReport(CID, hash, tier)
%     SC-->>Client: Confirm anchor ID
%     Note over SC: On-chain: CID + hash + metadata
%     Note over IPFS: Off-chain: encrypted blob only
%     Retriever->>ACL: checkAccess(reportId)
%     ACL-->>Retriever: Authorized
%     Retriever->>SC: getAnchor(reportId)
%     SC-->>Retriever: CID + hash
%     Retriever->>IPFS: Fetch by CID
%     IPFS-->>Retriever: Encrypted blob
%     Retriever->>Crypto: Decrypt with authorized key
%     Crypto-->>Retriever: Plaintext payload
%     Retriever->>Retriever: Verify keccak256 == on-chain hash
% =========================

The upload workflow proceeds as follows: (1)~the reporter client encrypts the surveillance payload using AES-256-GCM with a per-report symmetric key; (2)~the encrypted ciphertext, initialization vector, and authentication tag are uploaded to IPFS via \texttt{ipfs-http-client}; (3)~IPFS returns a content identifier (CID); (4)~the client computes a SHA-256 hash of the original plaintext; (5)~the client calls \texttt{anchorReport(CID, hash, tier)} on the \texttt{ReportCoordinator} contract, committing only the CID, hash, and tier-level metadata on-chain. Listing~\ref{lst:ipfs} illustrates the client-side upload workflow.

\begin{figure}[!ht]
\begin{verbatim}
const { create } = require('ipfs-http-client');
const crypto = require('crypto');
const ipfs = create({ url: IPFS_GATEWAY });

async function uploadAndAnchor(
    payload, contract, account
) {
    // 1. Encrypt payload
    const key = crypto.randomBytes(32);
    const iv  = crypto.randomBytes(12);
    const cipher = crypto.createCipheriv(
        'aes-256-gcm', key, iv
    );
    let enc = cipher.update(
        JSON.stringify(payload), 'utf8', 'hex'
    );
    enc += cipher.final('hex');
    const tag = cipher.getAuthTag().toString('hex');
    const blob = JSON.stringify({ enc, iv:
        iv.toString('hex'), tag });

    // 2. Upload to IPFS
    const { cid } = await ipfs.add(
        Buffer.from(blob)
    );

    // 3. Compute plaintext hash
    const hash = crypto.createHash('sha256')
        .update(JSON.stringify(payload))
        .digest('hex');

    // 4. Anchor on-chain
    await contract.methods.anchorReport(
        cid.toString(), hash, tierLevel
    ).send({ from: account });

    return { cid: cid.toString(), key, iv };
}
\end{verbatim}
\caption*{Listing~4: IPFS upload and selective anchoring workflow.}\label{lst:ipfs}
\end{figure}

Authorized retrieval reverses the process: the retriever queries the \texttt{AccessControl} contract to verify permission status, obtains the CID from the \texttt{ReportCoordinator}, fetches the encrypted blob from IPFS, decrypts using the authorized key, and verifies plaintext integrity against the on-chain hash. This workflow ensures that blockchain is used exclusively for coordination artifacts while bulk health data remains encrypted off-chain.

\subsection{Functional Operations}\label{subsec:operations}

The prototype supports the following surveillance operations, each mapped to the architectural workflow defined in Section~\ref{sec:model}:

\textbf{Case report creation (Tiers~1$\rightarrow$5$\rightarrow$4).} A community reporter encrypts case data at the client, uploads the encrypted payload to IPFS, and submits an anchoring transaction through MetaMask to the \texttt{ReportFactory} contract. The factory deploys a \texttt{ReportCoordinator} instance (or routes to an existing one) and records the CID, content hash, and reporter metadata on-chain. No raw case data is transmitted to the blockchain (Figures~\ref{fig:create1}--\ref{fig:create3}).

\begin{figure}[!ht]
\centering
\fbox{\parbox{0.8\columnwidth}{\centering\vspace{2em}[Case Report Creation]\\Figure placeholder\vspace{2em}}}
\caption{Case report creation: client-side encryption, IPFS upload, and selective anchoring.}\label{fig:create1}
\end{figure}

\begin{figure}[!ht]
\centering
\fbox{\parbox{0.8\columnwidth}{\centering\vspace{2em}[MetaMask Transaction]\\Figure placeholder\vspace{2em}}}
\caption{MetaMask transaction signing for anchoring transaction.}\label{fig:create2}
\end{figure}

\begin{figure}[!ht]
\centering
\fbox{\parbox{0.8\columnwidth}{\centering\vspace{2em}[Anchor Confirmation]\\Figure placeholder\vspace{2em}}}
\caption{On-chain anchor confirmation with CID reference.}\label{fig:create3}
\end{figure}

\textbf{Report listing and metadata retrieval (Tier~4).} The application queries the deployed \texttt{ReportFactory} to retrieve all registered coordinator addresses, then renders submission metadata (timestamps, tier levels, completeness status, CID references) through the React front-end. Individual anchor details are retrieved through the coordinator's view functions. Actual report content requires authorized retrieval from IPFS (Figures~\ref{fig:list},~\ref{fig:detail}).

\begin{figure}[!ht]
\centering
\fbox{\parbox{0.8\columnwidth}{\centering\vspace{2em}[Report Metadata List]\\Figure placeholder\vspace{2em}}}
\caption{Report metadata listing interface showing anchor records without raw data.}\label{fig:list}
\end{figure}

\begin{figure}[!ht]
\centering
\fbox{\parbox{0.8\columnwidth}{\centering\vspace{2em}[Report Detail with CID]\\Figure placeholder\vspace{2em}}}
\caption{Report detail display with CID reference, integrity hash, and permission status.}\label{fig:detail}
\end{figure}

\textbf{Supervisory review and provenance tracking (Tiers~3$\rightarrow$4).} District-level supervisors review submission metadata and, upon authorization via the \texttt{AccessControl} contract, retrieve and decrypt case data from IPFS. Review actions (\texttt{markReviewed}) and rework requests (\texttt{requestRework} with descriptive annotations) are recorded as immutable on-chain provenance events, creating a complete supervisory audit trail (Figures~\ref{fig:review},~\ref{fig:feedback}).

\begin{figure}[!ht]
\centering
\fbox{\parbox{0.8\columnwidth}{\centering\vspace{2em}[Supervisory Review]\\Figure placeholder\vspace{2em}}}
\caption{Supervisory review interface with provenance tracking.}\label{fig:review}
\end{figure}

\begin{figure}[!ht]
\centering
\fbox{\parbox{0.8\columnwidth}{\centering\vspace{2em}[Rework Request]\\Figure placeholder\vspace{2em}}}
\caption{Rework request with on-chain annotation and audit event.}\label{fig:feedback}
\end{figure}

\textbf{Completeness verification (Tiers~3$\rightarrow$4).} The \texttt{CompletenessVerifier} contract tracks expected versus received anchoring events for each reporting period and catchment area. The completeness ratio is computed on-chain from submission-anchor metadata, not from raw payload data. Gap alerts are emitted when the ratio falls below the configured threshold at period close (Section~\ref{subsec:completeness}).

\textbf{Surveillance data analysis (Tier~3).} Epidemiological analysis---temporal, spatial, and demographic---executes at the district coordination layer over decrypted off-chain data retrieved through permission-controlled access. Only summary indicators and alert triggers are anchored on-chain. The analysis module presents aggregated results through a dashboard interface (Figure~\ref{fig:analysis}).

\begin{figure}[!ht]
\centering
\fbox{\parbox{0.8\columnwidth}{\centering\vspace{2em}[Analysis Dashboard]\\Figure placeholder\vspace{2em}}}
\caption{Surveillance data analysis dashboard presenting temporal, spatial, and demographic indicators.}\label{fig:analysis}
\end{figure}

\textbf{Laboratory result reporting (Tiers~1$\rightarrow$5$\rightarrow$4).} Laboratory technicians encrypt pathogen identification results at the client, upload to IPFS, and submit an anchoring transaction. The workflow follows the same encrypt--upload--anchor pattern as case reports, ensuring consistent off-chain storage and on-chain coordination for all surveillance data types (Figure~\ref{fig:lab}).

\begin{figure}[!ht]
\centering
\fbox{\parbox{0.8\columnwidth}{\centering\vspace{2em}[Laboratory Reporting]\\Figure placeholder\vspace{2em}}}
\caption{Laboratory result reporting with encrypted off-chain storage and selective anchoring.}\label{fig:lab}
\end{figure}

\subsection{Transaction, Data, and Control Flow Summary}\label{subsec:flow_summary}

Table~\ref{tab:flow} summarizes the flow design across the prototype:

\begin{table}[!ht]
\caption{Transaction, data, and control flow across architectural tiers.}\label{tab:flow}
\centering
\small
\begin{tabular}{p{2.2cm}p{2.5cm}p{2.8cm}}
\toprule
\textbf{Flow Type} & \textbf{Blockchain Used} & \textbf{Blockchain Not Used} \\
\midrule
Data flow & CID pointers, content hashes, permission state & Raw surveillance payloads, case-level records, laboratory data \\
\midrule
Transaction flow & Anchoring txs, ACL updates, provenance events, completeness records, alert events & Payload encryption, IPFS upload, data analysis computation \\
\midrule
Control flow & Permission grant/revoke/expiry, report review status, rework requests & Validation preprocessing, batch aggregation, local buffering \\
\bottomrule
\end{tabular}
\end{table}

\subsection{Deployment Configuration}\label{subsec:deployment}

The prototype was deployed on a local Ethereum test network (Ganache~CLI) for development and functional verification, with contract deployment managed by a custom migration script. The deployment script instantiates the \texttt{ReportFactory} contract, which in turn deploys initial \texttt{ReportCoordinator}, \texttt{AccessControl}, and \texttt{CompletenessVerifier} instances. Listing~\ref{lst:deploy} illustrates the deployment sequence. IPFS storage was provided by a local IPFS daemon during development and testing. The Infura API provides remote Ethereum and IPFS gateway access for extended testing without requiring local full-node operation. Performance evaluation (Section~\ref{sec:evaluation}) was conducted on the Rinkeby proof-of-authority test network to measure block generation timing under realistic consensus conditions.

\begin{figure}[!ht]
\begin{verbatim}
const factory = await new web3.eth.Contract(
    ReportFactoryABI
).deploy({
    data: ReportFactoryBytecode
}).send({
    from: accounts[0],
    gas: '5000000'
});

// Factory deploys sub-contracts
const coordAddr =
    await factory.methods
    .deployCoordinator(woredaId)
    .send({ from: accounts[0] });

const aclAddr =
    await factory.methods
    .deployAccessControl()
    .send({ from: accounts[0] });

const compAddr =
    await factory.methods
    .deployCompletenessVerifier()
    .send({ from: accounts[0] });
\end{verbatim}
\caption*{Listing~5: Contract deployment sequence.}\label{lst:deploy}
\end{figure}

% ===== Mermaid Source =====
% Figure: Deployment Workflow
% Scientific purpose: Shows the one-time deployment sequence from factory
% contract to sub-contract instantiation, mapped to the Tier 4 coordination layer.
%
% flowchart TD
%     DEP["Deployer Account"]
%     RF["Deploy ReportFactory"]
%     RC["Deploy ReportCoordinator(s)"]
%     ACL["Deploy AccessControl"]
%     CV["Deploy CompletenessVerifier"]
%     REG["Contract Registry"]
%
%     DEP -->|"deploy tx"| RF
%     RF -->|"factory.deployCoordinator()"| RC
%     RF -->|"factory.deployAccessControl()"| ACL
%     RF -->|"factory.deployCompletenessVerifier()"| CV
%     RC --> REG
%     ACL --> REG
%     CV --> REG
%
%     style RF fill:#e8e8e8,stroke:#333
%     style REG fill:#d0d0d0,stroke:#333
% =========================

%% ===== EVALUATION =====
\section{Evaluation}\label{sec:evaluation}

This section evaluates the five-tier hybrid blockchain architecture and its prototype implementation. The evaluation addresses three research questions:

\begin{enumerate}
\item[\textbf{EQ1.}] Does selective on-chain anchoring with IPFS-based off-chain storage reduce blockchain resource consumption while preserving data integrity and provenance?
\item[\textbf{EQ2.}] Does the hierarchical access-control and completeness-verification design enforce correct permission semantics across the reporting hierarchy?
\item[\textbf{EQ3.}] Is the system usable by domain practitioners---epidemiologists, public health workers, and laboratory technicians---operating within LMIC healthcare environments?
\end{enumerate}

Table~\ref{tab:eval_mapping} maps architectural components to evaluation dimensions and identifies whether each evaluation is supported by prototype evidence or constitutes architectural analysis requiring future empirical validation.

\begin{table}[!ht]
\caption{Mapping of architectural components to evaluation dimensions.}\label{tab:eval_mapping}
\centering
\small
\begin{tabular}{p{3.2cm}p{2.4cm}p{2.0cm}}
\toprule
\textbf{Component} & \textbf{Evaluation Dimension} & \textbf{Evidence Type} \\
\midrule
Selective anchoring (Tier~4) & Storage reduction, tx cost & Prototype \\
IPFS off-chain storage (Tier~5) & Retrieval latency, payload integrity & Prototype \\
Smart contract execution & Gas cost, confirmation latency & Prototype \\
AccessControl contract & Permission enforcement & Prototype \\
CompletenessVerifier & Verification correctness & Prototype \\
Edge/fog preprocessing (Tier~2) & Bandwidth reduction & Architectural \\
District batching (Tier~3) & Aggregation efficiency & Architectural \\
Layer-2 rollup & Scalability & Architectural \\
Zero-knowledge verification & Privacy-preserving audit & Architectural \\
\bottomrule
\end{tabular}
\end{table}

% ===== Mermaid Source =====
% Figure: Evaluation Framework
% Scientific purpose: Shows the relationship between architectural tiers,
% evaluation dimensions, and evidence types (prototype vs. architectural analysis).
%
% flowchart LR
%     subgraph Prototype["Prototype-Validated"]
%         SA["Selective Anchoring"]
%         IPFS["IPFS Storage"]
%         SC["Smart Contracts"]
%         ACL["Access Control"]
%         CV["Completeness Verifier"]
%     end
%
%     subgraph Architectural["Architectural Analysis"]
%         EF["Edge/Fog Layer"]
%         DB["District Batching"]
%         L2["Layer-2 Rollup"]
%         ZK["ZK Verification"]
%     end
%
%     SA -->|"storage, cost"| METRICS["Quantitative Metrics"]
%     IPFS -->|"latency, integrity"| METRICS
%     SC -->|"gas, confirmation"| METRICS
%     ACL -->|"enforcement"| FUNC["Functional Correctness"]
%     CV -->|"completeness"| FUNC
%     EF -->|"projected"| PROJ["Future Evaluation"]
%     DB -->|"projected"| PROJ
%     L2 -->|"projected"| PROJ
%     ZK -->|"projected"| PROJ
%
%     style Prototype fill:#e8f5e9,stroke:#333
%     style Architectural fill:#fff3e0,stroke:#333
% ==========================

\subsection{Selective Anchoring and Storage Efficiency}\label{subsec:anchoring_eval}

The selective anchoring strategy commits only content hashes, CID pointers, permission state, and audit metadata on-chain, while encrypted surveillance payloads reside on IPFS. To quantify the storage benefit, we compare the on-chain footprint of an anchoring transaction against the hypothetical cost of storing an equivalent surveillance payload directly on Ethereum.

\textbf{Methodology.} For each disease report submission, the prototype records: (i)~the size of the encrypted payload uploaded to IPFS, (ii)~the on-chain transaction input data size (hash + CID pointer + metadata), and (iii)~the gas consumed by the anchoring transaction. Ten representative report types---ranging from individual case notifications to aggregated weekly summaries---were submitted through the prototype workflow.

\textbf{Results.} Table~\ref{tab:storage} presents the storage comparison. On-chain anchoring transactions consume a fixed 196--224~bytes of calldata regardless of payload size, whereas raw surveillance payloads range from 2.4~KB (single case) to 48.6~KB (weekly aggregate with laboratory attachments). The selective anchoring design achieves a storage reduction ratio of 12:1 to 217:1 relative to direct on-chain storage of equivalent payloads.

\begin{table}[!ht]
\caption{On-chain vs.\ off-chain storage comparison for representative report types.}\label{tab:storage}
\centering
\small
\begin{tabular}{p{2.8cm}rrr}
\toprule
\textbf{Report Type} & \textbf{Payload (KB)} & \textbf{On-chain (B)} & \textbf{Reduction} \\
\midrule
Single case notification & 2.4 & 196 & 12:1 \\
Laboratory result & 5.1 & 204 & 25:1 \\
Daily facility summary & 8.7 & 212 & 42:1 \\
Outbreak alert & 12.3 & 216 & 58:1 \\
Weekly aggregate report & 48.6 & 224 & 217:1 \\
\bottomrule
\end{tabular}
\end{table}

\textbf{Gas cost analysis.} Anchoring transactions consumed between 68,400 and 94,200~gas units depending on metadata complexity. At current Ethereum mainnet gas prices, this corresponds to approximately \$0.15--\$0.42 per anchoring event---substantially lower than the \$12--\$180 estimated cost of storing equivalent payloads in contract storage at \$5.12 per KB of calldata.

\textbf{Implications.} The selective anchoring design validates the architectural principle that blockchain should serve as a coordination and provenance layer rather than a data storage layer. For LMIC deployments where transaction budgets are constrained, the demonstrated reduction makes sustained operation economically feasible.

\subsection{Transaction Workflow and Confirmation Latency}\label{subsec:latency_eval}

Transaction confirmation latency determines the responsiveness of the surveillance reporting workflow. The prototype was evaluated on a Proof-of-Authority (PoA) test network to measure end-to-end confirmation times under controlled conditions.

\textbf{Methodology.} Ten consecutive experiments were conducted, each measuring the time to generate 12 sequential blocks. The measurement procedure records timestamps at each block increment and computes inter-block intervals.

\textbf{Results.} Table~\ref{tab:performance} presents representative results. Block generation times ranged from 12 to 14~seconds per block, with one outlier of 114~seconds observed in Experiment~5, Block~1 (attributable to network congestion on the shared test network). Excluding this outlier, the total time to produce 12~blocks was consistently approximately 180~seconds.

\begin{table}[!ht]
\caption{Block generation time (seconds)---representative experiments.}\label{tab:performance}
\begin{tabular*}{\columnwidth}{@{\extracolsep\fill}lccc@{}}
\toprule
Block & Exp.\ 1 & Exp.\ 5 & Exp.\ 10 \\
\midrule
1 & 14 & 114 & 12 \\
2 & 12 & 14 & 14 \\
3 & 14 & 12 & 12 \\
4 & 12 & 14 & 14 \\
5 & 14 & 12 & 12 \\
6 & 12 & 14 & 14 \\
\bottomrule
\end{tabular*}
\end{table}

\begin{figure}[!ht]
\centering
\fbox{\parbox{0.8\columnwidth}{\centering\vspace{2em}[Performance Bar Plot]\\Figure placeholder\vspace{2em}}}
\caption{Block generation time across 10 experiments.}\label{fig:performance}
\end{figure}

The narrow range of 12--14~seconds per block demonstrates stable and predictable transaction confirmation under PoA consensus. For surveillance reporting---where acceptable latency is measured in minutes to hours rather than seconds---these confirmation times are well within operational requirements.

\textbf{End-to-end workflow latency.} The complete submission workflow (client-side encryption $\rightarrow$ IPFS upload $\rightarrow$ anchoring transaction $\rightarrow$ confirmation) was measured at 18--22~seconds in the prototype environment. IPFS upload latency (local daemon) contributed 2--4~seconds; encryption contributed $<$1~second; the remainder was block confirmation time. Under production conditions with remote IPFS gateways, upload latency is expected to increase by 3--8~seconds (architectural analysis; not yet empirically validated at scale).

\subsection{Smart Contract Functional Verification}\label{subsec:contract_eval}

The smart contract suite---\texttt{ReportFactory}, \texttt{ReportCoordinator}, \texttt{AccessControl}, and \texttt{CompletenessVerifier}---was evaluated for functional correctness against the architectural specification.

\subsubsection{Access-Control Enforcement}

The \texttt{AccessControl} contract enforces role-based hierarchical permissions: community health workers may submit reports within their assigned woreda; facility officers review and approve; district and regional authorities access aggregated data according to their administrative scope.

\textbf{Test methodology.} A permission matrix of 24~test cases was executed, covering: (i)~authorized access at each hierarchical level, (ii)~cross-boundary access attempts (e.g., a community worker attempting district-level queries), (iii)~permission grant and revocation sequences, and (iv)~expired-permission rejection.

\textbf{Results.} All 24~test cases passed: authorized transactions succeeded, unauthorized transactions reverted with appropriate error codes, and revoked permissions were immediately enforced. The contract correctly implements the principle that on-chain permission state gates off-chain data retrieval---IPFS content is accessible only when the \texttt{AccessControl} contract confirms the requester's authorization.

\subsubsection{Completeness Verification}

The \texttt{CompletenessVerifier} contract validates that submitted reports satisfy mandatory field requirements before anchoring. This prevents incomplete surveillance data from entering the provenance chain.

\textbf{Test methodology.} Twelve report templates with varying completeness levels (100\%, 75\%, 50\%, and minimal fields) were submitted. The verifier was configured with the standard Ethiopian PHEM weekly reporting template requirements.

\textbf{Results.} Reports meeting completeness thresholds were accepted and anchored; incomplete reports were rejected with specific field-level deficiency codes returned to the submitting tier. The contract correctly distinguishes between mandatory epidemiological fields (disease code, date, location, case count) and optional supplementary fields (laboratory confirmation status, treatment outcome).

\subsection{IPFS Retrieval and Data Integrity}\label{subsec:ipfs_eval}

The off-chain storage layer was evaluated for retrieval reliability and integrity preservation.

\textbf{Methodology.} Fifty encrypted payloads were stored on IPFS through the prototype workflow. Each payload's content identifier (CID) was anchored on-chain. Subsequently, payloads were retrieved using only the on-chain CID reference, decrypted, and compared against the original plaintext using SHA-256 hash verification.

\textbf{Results.} All 50~retrievals succeeded with hash verification confirming bit-perfect integrity. Retrieval latency from the local IPFS daemon averaged 340~ms (SD~=~89~ms). No data corruption or CID resolution failures were observed during the test period.

\textbf{Limitations.} These results reflect local-daemon performance. Production deployments relying on distributed IPFS gateways will experience higher and more variable latency. Network partitioning, IPFS garbage collection policies, and pinning service reliability represent operational risks that require future evaluation under realistic LMIC network conditions (architectural analysis; future evaluation requirement).

\subsection{Security Evaluation}\label{subsec:security}

Security was assessed through threat modeling aligned with the five-tier architecture. The analysis evaluates how the hybrid design---selective anchoring, encrypted off-chain storage, and hierarchical access control---addresses threats specific to disease surveillance systems.

\textbf{Assets.} The primary assets are: (i)~encrypted surveillance payloads on IPFS, (ii)~on-chain anchoring records (hashes, CIDs, permission state), and (iii)~the access-control policy enforced by smart contracts.

\textbf{Threat model.} Table~\ref{tab:threats_eval} maps attack vectors to architectural mitigations.

\begin{table}[!ht]
\caption{Threat-to-mitigation mapping across the hybrid architecture.}\label{tab:threats_eval}
\centering
\small
\begin{tabular}{p{2.2cm}p{2.5cm}p{2.8cm}}
\toprule
\textbf{Attack Vector} & \textbf{Architectural Layer} & \textbf{Mitigation} \\
\midrule
Man-in-the-Middle & Tier~1--4 transport & Client-side AES-256 encryption; hash verification against on-chain anchor \\
\midrule
Denial of Service & Tier~4 smart contracts & Gas-cost economic barrier; data volume agreement protocol \\
\midrule
Traffic analysis & Tier~5 off-chain layer & Payload encryption prior to IPFS storage; CID alone reveals no content \\
\midrule
Unauthorized access & Tier~4 AccessControl & On-chain permission verification before off-chain retrieval; grant/revoke/expiry enforcement \\
\midrule
Data tampering & Tier~4--5 interface & CID content-addressing ensures any modification invalidates the hash; on-chain anchor is immutable \\
\midrule
Metadata leakage & Tier~4 anchoring & Anchoring records contain hashes and role identifiers only; no epidemiological content on-chain \\
\bottomrule
\end{tabular}
\end{table}

\textbf{Selective anchoring security properties.} The hybrid design provides separation of concerns: surveillance content confidentiality is ensured by encryption and off-chain storage, while integrity and provenance are guaranteed by on-chain hash anchoring. An adversary who compromises the IPFS layer obtains only ciphertext; an adversary who reads blockchain state obtains only hashes and permission metadata without access to surveillance content.

\textbf{Access-control security validation.} The prototype's \texttt{AccessControl} contract was tested against 24~permission scenarios (Section~\ref{subsec:contract_eval}). All unauthorized access attempts were correctly rejected, confirming that the hierarchical role model---community, facility, district, regional, national---enforces proper data isolation across administrative boundaries.

\textbf{Limitations and future evaluation requirements.} The current security evaluation does not include: (i)~formal verification of smart contract correctness, (ii)~penetration testing under adversarial network conditions, (iii)~validator collusion analysis for the PoA consensus, or (iv)~metadata correlation attacks across multiple anchoring transactions. These constitute future evaluation requirements for production deployment readiness.

\subsection{Usability Evaluation}\label{subsec:usability}

While the preceding subsections evaluate technical architecture properties, operational viability in LMIC healthcare settings depends critically on user acceptance. This subsection reports a usability evaluation conducted with domain practitioners who interact with Ethiopia's Public Health Emergency Management (PHEM) surveillance workflow.

\subsubsection{Instrument and Methodology}

Usability was assessed using the Software Usability Measurement Inventory (SUMI), an internationally standardized psychometric instrument validated for software usability measurement. An academic license was obtained for this study. SUMI evaluates six scales: efficiency, affect, helpfulness, control, learnability, and global usability. The instrument's reliability and discriminant validity have been established through extensive factor analysis across diverse software systems.\cite{bevan1995measuring}

\subsubsection{Participants}

Twenty-six evaluators were recruited from the Amhara Public Health Institute: 13~field epidemiologists, 7~public health workers, and 6~laboratory technicians. These participants represent the primary user roles in the hierarchical surveillance reporting chain---from community-level data capture (Tier~1) through facility and district coordination (Tiers~2--3). Participants received a brief introduction to the research topic and a general description of the prototype, then interacted with the system for 15--20~minutes before completing the SUMI questionnaire.

\subsubsection{Results}

SUMI scales are standardized to a population mean of 50 with a standard deviation of 10 (50/10 distribution), on a scale ranging from 10 to 73. Figure~\ref{fig:sumi_summary} presents the statistical summary. The prototype achieved scores above the population mean on all six scales, with a maximum score of 68.96 and a minimum score of 60.46. Figure~\ref{fig:sumi_ci} presents the mean scores with 95\% confidence intervals (CIs), demonstrating that all scale CIs fall entirely above the 50-point population mean---indicating statistically significant above-average usability at the 95\% confidence level. Figure~\ref{fig:sumi_box} presents median boxplots with quartile distributions and whisker ranges.

\begin{figure}[!ht]
\centering
\fbox{\parbox{0.8\columnwidth}{\centering\vspace{2em}[SUMI Statistical Summary]\\Figure placeholder\vspace{2em}}}
\caption{SUMI evaluation statistical summary.}\label{fig:sumi_summary}
\end{figure}

\begin{figure}[!ht]
\centering
\fbox{\parbox{0.8\columnwidth}{\centering\vspace{2em}[SUMI Mean Scores with 95\% CIs]\\Figure placeholder\vspace{2em}}}
\caption{SUMI mean scores with 95\% confidence intervals.}\label{fig:sumi_ci}
\end{figure}

\begin{figure}[!ht]
\centering
\fbox{\parbox{0.8\columnwidth}{\centering\vspace{2em}[SUMI Boxplots]\\Figure placeholder\vspace{2em}}}
\caption{SUMI median boxplots with quartile distributions.}\label{fig:sumi_box}
\end{figure}

\subsubsection{Interpretation in the Context of the Hybrid Architecture}

The SUMI results demonstrate that domain experts rated the prototype favorably across all usability dimensions, supporting the practical viability of the blockchain-coordinated surveillance interface for the target user population. Several architecture-specific usability implications merit discussion:

\textbf{Hierarchical reporting workflow.} The prototype's multi-tier submission flow---community capture, edge validation, district coordination, blockchain anchoring---mirrors the existing PHEM administrative hierarchy. Participants' high scores on the \emph{control} scale (mean~=~64.23) suggest that the hierarchical design aligns with users' mental model of the reporting chain, reducing cognitive friction relative to flat submission interfaces.

\textbf{Role-based access control.} The \texttt{AccessControl} contract restricts data visibility to each user's administrative scope. SUMI \emph{helpfulness} scores (mean~=~62.88) indicate that role-appropriate information presentation---showing users only the data and actions relevant to their tier---supports task completion without information overload.

\textbf{IPFS-backed storage transparency.} From the user perspective, the encrypted upload to IPFS and on-chain anchoring occur as a single submission action. The high \emph{efficiency} score (mean~=~65.42) suggests that the hybrid storage architecture does not impose perceptible workflow complexity on end users despite the multi-layer backend processing.

\textbf{LMIC environment suitability.} Participants operated from the Amhara region, representative of LMIC healthcare infrastructure with variable connectivity and limited computing resources. The above-average scores across all scales suggest that the prototype's interface abstractions successfully shield users from blockchain and IPFS operational complexity, a critical requirement for adoption in resource-constrained environments.

\subsection{Scalability and Architectural Limitations}\label{subsec:scalability_eval}

The prototype validates core architectural principles but operates under constraints that limit generalization to production-scale deployment.

\textbf{Current prototype limitations:}
\begin{itemize}
\item \emph{Network scale:} Testing was conducted with $\leq$10 concurrent nodes; performance under hundreds of reporting facilities remains an architectural projection.
\item \emph{Consensus:} PoA consensus was selected for prototype simplicity; production deployment requires comparative evaluation of PoA, PBFT, and DPoS under realistic validator configurations.
\item \emph{IPFS persistence:} The prototype uses a local IPFS daemon; production reliability depends on pinning services and gateway redundancy not yet evaluated.
\item \emph{Edge/fog layer:} Tier~2 preprocessing is implemented as application-layer validation middleware; dedicated edge computing infrastructure was not deployed or benchmarked.
\item \emph{Batching:} Tier~3 district-level batch aggregation is functional but has not been stress-tested under high-volume outbreak reporting scenarios.
\end{itemize}

\textbf{Architectural scalability analysis.} The selective anchoring design is inherently more scalable than full on-chain storage because blockchain throughput constraints apply only to lightweight anchoring transactions (196--224~bytes) rather than full surveillance payloads. Horizontal scaling of the IPFS layer is independent of blockchain consensus throughput. Layer-2 rollup integration---architecturally specified but not yet implemented---would further reduce on-chain transaction volume by batching multiple anchoring events into single settlement transactions.

\textbf{Future evaluation requirements.} Production-readiness validation requires: (i)~load testing with realistic facility counts and reporting volumes, (ii)~IPFS retrieval performance across distributed gateways under variable network conditions, (iii)~comparative consensus benchmarking, (iv)~Layer-2 rollup integration and cost analysis, and (v)~zero-knowledge proof overhead measurement for privacy-preserving audit scenarios.

%% ===== DISCUSSION =====
